\section{Exterior extensions and their boundary resource}
\label{jext:section}

Exterior addition has precedents in BBL and Blanc~\cite{bbl2008,blanc2011}.
The contribution here is an explicit geometric certificate, a quantitative
defect formulation for both parities, and an obstruction to repeatedly
demanding the full gain proposed in March~\cite{zarzuelo2026}.

Let $A$ be a simple $n$-line arrangement, $n\ge3$, with equations $f_i=0$ and
vertices $p_{ij}$.  A linear functional $w$ is admissible if it is nonzero
on every line direction.  At $p_{ij}$ choose support directions $u_i,u_j$
with $w(u_i)=w(u_j)=1$.  Call the pair \emph{visible} if, for every old
line $f_k=0$, the numbers
\[
 f_k(p_{ij}),\quad Df_k(u_i),\quad Df_k(u_j)
\]
have one common weak sign.  This condition certifies an empty unbounded
wedge using only old data.

\begin{theorem}[Verified exterior realization]
\label{jext:visible}
If $A$ contains $T$ distinct certified triangles and $c$ distinct pairs
visible from an admissible $w$, adjoining $w(x)=H$, where
$H>\max_{i<j}w(p_{ij})$, gives a simple arrangement with at least $T+c$
triangles.
\end{theorem}
\begin{proof}
The new line is nonparallel to every support, misses all old vertices,
and lies beyond every old bounded cell.  Thus simplicity and old triangles
are preserved.  A visible pair has new vertices
$p_{ij}+(H-w(p_{ij}))u_i$ and $p_{ij}+(H-w(p_{ij}))u_j$.
Their affine evaluations are
$f_k(p_{ij})+(H-w(p_{ij}))Df_k(u)$, so all three vertices have one weak
sign against every old line.  The triangle is nondegenerate and empty.
Distinct pairs give distinct new support triples, all containing the
new line and hence distinct from the old triples.
\end{proof}
This is \lean{Exterior.extension}, a real-coordinate Lean theorem with
no parity restriction or assumed future count.

Write $T=t(A)$, $d=n(n-2)-3T$, and let $b$ count unbounded two-sided cells.
The quantity $d$ counts unused bounded elementary segments: a segment
in a simple arrangement supports at most one triangle.

\begin{proposition}[Boundary-defect extension]
\label{jext:defect}
For $n\ge3$, define
\begin{equation}
 g(n,T)=\ceil{\frac{n-1}{2n}\max\{3,\,3T-n(n-3)\}}.
 \label{jext:gain}
\end{equation}
Then $\Ks(n)\ge T$ implies $\Ks(n+1)\ge T+g(n,T)$.
\end{proposition}
\begin{proof}
There are $2n$ free rays, of which $2b$ belong to wedges.  At an unpaired
ray's endpoint $L\cap M$, both incident segments of $M$ are bounded.
They cannot both support triangles, since those triangles would share
the unique bounded segment of $L$ starting there.  Charge this ray to
an unused segment of $M$.  Each unused segment receives at most two
endpoint charges, giving $b+d\ge n$.  Convex-hull vertices of the finite
intersection set supply at least three wedges, so $b\ge\max\{3,n-d\}$.

The two rays of a wedge are consecutive in cyclic direction order.
Among the $2n$ admissible normal sectors, each wedge is visible in
exactly $n-1$: positivity selects $n$ consecutive oriented directions.
Thus some sector sees at least $\ceil{(n-1)b/(2n)}$ wedges.  Apply
Theorem~\ref{jext:visible}.  The expression is nondecreasing in $T$,
so a lower count can replace the witness's actual count.
\end{proof}
The cyclic count and defect arithmetic are formalized.  The universal
ray-charging and geometric cyclic-data extraction remain manuscript
proofs; the proposition is not claimed as an end-to-end Lean theorem.

If $n=2m+1$ and $d\le2$, then $b\ge2m-1$.  Averaging forces at least
$m$ visible wedges, since
$2m(2m-1)>(4m+2)(m-1)$; no sector can contain more than $m$ disjoint
ray pairs.  Hence every such odd input admits the full gain $m$.
For either parity the exact full-gain test is that some sector contains
$\floor{n/2}$ wedge pairs.  The defect rule itself can be iterated
without assuming this stronger condition persists.

Repeated exterior additions have a finite boundary resource.  If $b_r$
is the wedge count and $c_r$ the gain, each new triangle consumes an
old wedge.  No new free ray appears at an old vertex, and new wedges
can use only the two free rays of the added line.  Therefore
\begin{equation}
 b_{r+1}+c_r\le b_r+2,\qquad
 \sum_{j<r}c_j\le b_0+2r-3.
 \label{jext:budget}
\end{equation}
Quadratically accumulating full gains cannot continue indefinitely.
For the saved $49:767$ seed, $b_0=47$; gains $24$ and $25$ in two
exterior steps would leave at most two wedges, contradicting $b\ge3$.
This does not refute a recurrence for maxima using reconstruction.

\begin{figure}[tb]
\centering
\includegraphics[width=.82\linewidth]{successor-49-resource.pdf}
\caption{One exact exterior chain from $49:767$. Its gains $24,23,2,2$ produce counts $791,814,816,818$. The displayed capacities belong to these particular arrangements; the general obstruction is the wedge budget~\eqref{jext:budget}, not an upper bound for $\K$.}\label{jext:figure}
\end{figure}

Conditionally, $F(n+1)\ge F(n)+\floor{n/2}$ yields
\[
 F(n)\ge F(s)+\floor{(n-1)^2/4}-\floor{(s-1)^2/4}.
\]
The recurrence remains unproved for $\K$.  Nevertheless its $49$-seed
numerical target is independently unconditional:
\begin{equation}
 \K(n)\ge\floor{(n-1)^2/4}+191\qquad(n\ge49).
 \label{jext:49}
\end{equation}
Lean proves this from the finite 49/50 witnesses and the classical
baseline for $n\ge51$; it does not construct a nested full-gain chain.
